\textbf{Phần 1. Trắc nghiệm lựa chọn}
\setcounterpageref{ex}{0}

\begin{ex} %[5D7N41]
Muốn tìm $30\%$ của số $80$, ta thực hiện
\choice
{$80 : 30 \times 100$}
{$80 \times 100 : 30$}
{\True $80 \times 30 : 100$}
{$80 + 30 : 100$}
\loigiai{
Muốn tìm $30\%$ của $80$, ta lấy $80 \times 30 : 100 = 24$.\\
\textbf{Đáp án của câu hỏi là C}
}
\end{ex}

\begin{ex} %[5D7N41]
$15\%$ của $240\text{ kg}$ là
\choice
{\True $36\text{ kg}$}
{$360\text{ kg}$}
{$3{,}6\text{ kg}$}
{$24\text{ kg}$}
\loigiai{
Ta có: $240 \times 15 : 100 = 36\text{ (kg)}$.\\
\textbf{Đáp án của câu hỏi là A}
}
\end{ex}

\begin{ex} %[5D7H41]
Một lớp có $35$ học sinh, trong đó có $60\%$ là học sinh nữ. Số học sinh nữ của lớp đó là
\choice
{$14$ bạn}
{$20$ bạn}
{$24$ bạn}
{\True $21$ bạn}
\loigiai{
Số học sinh nữ là: $35 \times 60 : 100 = 21$ (bạn).\\
\textbf{Đáp án của câu hỏi là D}
}
\end{ex}

\begin{ex} %[5D7H41]
Một chiếc áo có giá niêm yết $250\,000$ đồng, nhân dịp khuyến mại được giảm giá $20\%$. Số tiền được giảm giá là
\choice
{\True $50\,000$ đồng}
{$40\,000$ đồng}
{$60\,000$ đồng}
{$200\,000$ đồng}
\loigiai{
Số tiền giảm giá là: $250\,000 \times 20 : 100 = 50\,000$ (đồng).\\
\textbf{Đáp án của câu hỏi là A}
}
\end{ex}

\begin{ex} %[5D7H41]
Diện tích một khu vườn là $450\text{ m}^2$, người ta dùng $35\%$ diện tích để trồng hoa. Diện tích đất trồng hoa là
\choice
{$140\text{ m}^2$}
{\True $157{,}5\text{ m}^2$}
{$160\text{ m}^2$}
{$135\text{ m}^2$}
\loigiai{
Diện tích trồng hoa là: $450 \times 35 : 100 = 157{,}5\text{ (m}^2\text{)}$.\\
\textbf{Đáp án của câu hỏi là B}
}
\end{ex}

\begin{ex} %[5D7H41]
Một cửa hàng có $650\text{ kg}$ hoa quả. Trong ngày đầu tiên cửa hàng bán được $40\%$ số hoa quả đó. Khối lượng hoa quả còn lại sau ngày đầu là
\choice
{$260\text{ kg}$}
{$350\text{ kg}$}
{$400\text{ kg}$}
{\True $390\text{ kg}$}
\loigiai{
Số hoa quả đã bán: $650 \times 40 : 100 = 260\text{ (kg)}$.\\
Số hoa quả còn lại: $650 - 260 = 390\text{ (kg)}$.\\
\textbf{Đáp án của câu hỏi là D}
}
\end{ex}

\begin{ex} %[5D7V41]
Một mảnh đất hình chữ nhật có chiều dài $30\text{ m}$, chiều rộng bằng $60\%$ chiều dài. Chu vi mảnh đất đó là
\choice
{$48\text{ m}$}
{\True $96\text{ m}$}
{$72\text{ m}$}
{$540\text{ m}$}
\loigiai{
Chiều rộng mảnh đất: $30 \times 60 : 100 = 18\text{ (m)}$.\\
Chu vi mảnh đất: $(30 + 18) \times 2 = 96\text{ (m)}$.\\
\textbf{Đáp án của câu hỏi là B}
}
\end{ex}

\begin{ex} %[5D7V41]
Một đôi giày có giá niêm yết là $400\,000$ đồng. Cửa hàng giảm giá $15\%$. Khách hàng cần trả số tiền để mua đôi giày đó là
\choice
{$60\,000$ đồng}
{$360\,000$ đồng}
{\True $340\,000$ đồng}
{$350\,000$ đồng}
\loigiai{
Số tiền giảm giá: $400\,000 \times 15 : 100 = 60\,000$ (đồng).\\
Số tiền phải trả: $400\,000 - 60\,000 = 340\,000$ (đồng).\\
\textbf{Đáp án của câu hỏi là C}
}
\end{ex}

\begin{ex} %[5D7V41]
Bác Bình gửi tiết kiệm $80\,000\,000$ đồng tại ngân hàng với lãi suất $6\%$ một năm. Sau một năm, bác Bình nhận về tất cả bao nhiêu tiền cả gốc lẫn lãi?
\choice
{$4\,800\,000$ đồng}
{$86\,000\,000$ đồng}
{$88\,000\,000$ đồng}
{\True $84\,800\,000$ đồng}
\loigiai{
Tiền lãi sau một năm: $80\,000\,000 \times 6 : 100 = 4\,800\,000$ (đồng).\\
Tổng số tiền nhận được: $80\,000\,000 + 4\,800\,000 = 84\,800\,000$ (đồng).\\
\textbf{Đáp án của câu hỏi là D}
}
\end{ex}

\begin{ex} %[5D7V41]
Một cuốn sách dày $240$ trang. Ngày đầu Lan đọc được $25\%$ tổng số trang, ngày thứ hai Lan đọc được $40\%$ số trang còn lại. Số trang sách Lan chưa đọc là
\choice
{\True $108$ trang}
{$60$ trang}
{$72$ trang}
{$120$ trang}
\loigiai{
Ngày đầu đọc được: $240 \times 25 : 100 = 60$ (trang).\\
Số trang còn lại sau ngày đầu: $240 - 60 = 180$ (trang).\\
Ngày thứ hai đọc được: $180 \times 40 : 100 = 72$ (trang).\\
Số trang chưa đọc: $180 - 72 = 108$ (trang).\\
\textbf{Đáp án của câu hỏi là A}
}
\end{ex}

\begin{ex} %[5D7V41]
Một bể bơi chứa $120\text{ m}^3$ nước, chiếm $75\%$ dung tích của bể. Dung tích toàn bộ bể bơi đó là
\choice
{$90\text{ m}^3$}
{$150\text{ m}^3$}
{\True $160\text{ m}^3$}
{$180\text{ m}^3$}
\loigiai{
Dung tích bể bơi là: $120 : 75 \times 100 = 160\text{ (m}^3\text{)}$.\\
\textbf{Đáp án của câu hỏi là C}
}
\end{ex}

\begin{ex} %[5D7C41]
Giá một chiếc tivi là $12\,000\,000$ đồng. Nhân dịp khai trương, siêu thị giảm giá $10\%$. Sau đó nhân dịp Tết, siêu thị lại giảm thêm $5\%$ trên giá đã giảm. Giá chiếc tivi sau hai lần giảm là
\choice
{$10\,200\,000$ đồng}
{\True $10\,260\,000$ đồng}
{$10\,500\,000$ đồng}
{$10\,800\,000$ đồng}
\loigiai{
Giá sau lần giảm thứ nhất: $12\,000\,000 \times (100\% - 10\%) = 10\,800\,000$ (đồng).\\
Giá sau lần giảm thứ hai: $10\,800\,000 \times (100\% - 5\%) = 10\,260\,000$ (đồng).\\
\textbf{Đáp án của câu hỏi là B}
}
\end{ex}

\textbf{Phần 2. Trắc nghiệm đúng sai}
\setcounterpageref{ex}{0}

\begin{ex} %[5D7H41]
Một cửa hàng nhập về $600\text{ kg}$ gạo. Buổi sáng bán được $30\%$ số gạo đó, buổi chiều bán được $45\%$ số gạo đó. Xét tính đúng sai:
\choiceTF
{\True Buổi sáng bán được $180\text{ kg}$ gạo.}
{Buổi chiều bán được $250\text{ kg}$ gạo.}
{\True Cả hai buổi bán được $450\text{ kg}$ gạo.}
{\True Số gạo còn lại chiếm $25\%$ tổng số gạo ban đầu.}
\loigiai{
\begin{itemize}
    \item Buổi sáng: $600 \times 30\% = 180\text{ kg}$. Đúng.
    \item Buổi chiều: $600 \times 45\% = 270\text{ kg} \neq 250\text{ kg}$. Sai.
    \item Cả hai buổi: $180 + 270 = 450\text{ kg}$. Đúng.
    \item Tỉ lệ còn lại: $100\% - (30\% + 45\%) = 25\%$. Đúng.
\end{itemize}
\begin{itemize}
    \item a) \textbf{Đúng}.
    \item b) \textbf{Sai}.
    \item c) \textbf{Đúng}.
    \item d) \textbf{Đúng}.
\end{itemize}
}
\end{ex}

\begin{ex} %[5D7V41]
Một trường tiểu học có $800$ học sinh. Trong đó có $48\%$ là học sinh nam, số học sinh nữ đạt danh hiệu học sinh xuất sắc chiếm $25\%$ tổng số học sinh nữ. Xét tính đúng sai:
\choiceTF
{\True Trường đó có $384$ học sinh nam.}
{\True Trường đó có $416$ học sinh nữ.}
{Số học sinh nữ xuất sắc là $100$ bạn.}
{\True Số học sinh nữ xuất sắc là $104$ bạn.}
\loigiai{
\begin{itemize}
    \item Học sinh nam: $800 \times 48\% = 384$ bạn. Đúng.
    \item Học sinh nữ: $800 - 384 = 416$ bạn. Đúng.
    \item Nữ xuất sắc: $416 \times 25\% = 104$ bạn $\neq 100$ bạn. Ý 3 Sai, Ý 4 Đúng.
\end{itemize}
\begin{itemize}
    \item a) \textbf{Đúng}.
    \item b) \textbf{Đúng}.
    \item c) \textbf{Sai}.
    \item d) \textbf{Đúng}.
\end{itemize}
}
\end{ex}

\textbf{Phần 3. Trả lời ngắn}
\setcounterpageref{ex}{0}

\begin{ex} %[5D7H41]
Tìm $24\%$ của $350$.
\shortans{84}
\loigiai{
\textbf{Đáp số:} $84$\\
$350 \times 24 : 100 = 84$.
}
\end{ex}

\begin{ex} %[5D7H41]
Một mảnh đất có diện tích $520\text{ m}^2$, người ta dùng $45\%$ diện tích để trồng rau. Tính diện tích đất trồng rau theo mét vuông.
\shortans{234}
\loigiai{
\textbf{Đáp số:} $234$\\
$520 \times 45 : 100 = 234\text{ m}^2$.
}
\end{ex}

\begin{ex} %[5D7V41]
Một chiếc xe máy có giá niêm yết $28\,000\,000$ đồng, nhân dịp lễ được giảm giá $5\%$. Hỏi khách hàng phải trả bao nhiêu triệu đồng để mua chiếc xe máy đó? (Nhập số thập phân có dấu phẩy, ví dụ: 26,6).
\shortans{26,6}
\loigiai{
\textbf{Đáp số:} $26{,}6$\\
$28 \times (100 - 5) : 100 = 26{,}6\text{ triệu đồng}$.
}
\end{ex}

\begin{ex} %[5D7V41]
Một hình chữ nhật có chiều dài $40\text{ m}$, chiều rộng bằng $70\%$ chiều dài. Tính diện tích hình chữ nhật đó theo mét vuông.
\shortans{1120}
\loigiai{
\textbf{Đáp số:} $1120$\\
Chiều rộng: $40 \times 70\% = 28\text{ m}$.\\
Diện tích: $40 \times 28 = 1\,120\text{ m}^2$.
}
\end{ex}

\begin{ex} %[5D7V41]
Biết $35\%$ của số $x$ là $147$. Tìm số $x$.
\shortans{420}
\loigiai{
\textbf{Đáp số:} $420$\\
$x = 147 : 35 \times 100 = 420$.
}
\end{ex}

\begin{ex} %[5D7C41]
Một xí nghiệp có $300$ công nhân. Sau một năm, số công nhân tăng thêm $8\%$. Sang năm tiếp theo, số công nhân lại tăng thêm $5\%$ so với năm trước đó. Hỏi sau hai năm xí nghiệp đó có tất cả bao nhiêu công nhân?
\shortans{340 hoặc 340,2 (khoảng 340 người)}
\loigiai{
\textbf{Đáp số:} $340$ hoặc $340{,}2$ (khoảng 340 người)\\
Sau năm 1: $300 \times 1{,}08 = 324$ người.\\
Sau năm 2: $324 \times 1{,}05 = 340{,}2 \approx 340$ người.
}
\end{ex}
